\subsection{The Fourier transform}

We keep the notation of the previous number.

Let \(\pi \in \widehat{\mathbf{G}}\) . The vector space \(\operatorname{End}(\mathrm{E}_{\pi})\) is equipped with the structure of a Hilbert space whose scalar product is given by

\[
\langle u _ {1} \mid u _ {2} \rangle = \dim (\pi) \operatorname{Tr} (u _ {1} ^ {*} u _ {2}) = \dim (\pi) \operatorname{Tr} (u _ {2} u _ {1} ^ {*})
\]

for \(u_{1}\) , \(u_{2}\) in \(\operatorname{End}(\mathrm{E}_{\pi})\) (cf. EVT, V, p. 52, th. 1).

We denote by \(\|u\|_{2} = \sqrt{\langle u|u\rangle}\) the norm of an element u of \(\operatorname{End}(\operatorname{E}_{\pi})\) for \(\pi \in \widehat{G}\) . For every \(g \in G\) , we have \(\|\pi(g)\|_{2} = \dim(\pi)\) since \(\pi(g)\) is unitary.

The norm denoted here \(\|u\|_{2}\) differs by a factor \(\dim(\pi)\) from the norm defined in EVT, V, p. 52 on the space of Hilbert–Schmidt operators on \(E_{\pi}\) .

Lemma 4. --- Let \(\pi\) be an irreducible representation of G. The map \(f \mapsto \pi(f)\) defines, by passing to subspaces, an isometric isomorphism from \(\varrho_{\mathrm{G}}(\overline{\pi})\) onto \(\operatorname{End}(\operatorname{E}_{\pi})\) .

Set \(\varepsilon_{\pi}(g,h)u = \pi(g)\circ u\circ\pi(h^{-1})\) for all \((g,h)\in\mathrm{G}\times\mathrm{G}\) and for every \(u\in\mathrm{End}(\mathrm{E}_{\pi})\) . The map \(\varepsilon_{\pi}\) is a continuous representation of \(\mathrm{G}\times\mathrm{G}\) in \(\mathrm{End}(\mathrm{E}_{\pi})\) . It is unitary. Indeed, since \(\pi\) itself is unitary, it follows that

\[
\begin{array}{c} \| \varepsilon_ {\pi} (g, h) u \| _ {2} ^ {2} = \dim (\pi) \operatorname{Tr} \Big ((\pi (g) u \pi (h ^ {- 1})) ^ {*} \pi (g) u \pi (h ^ {- 1}) \Big) \\ = \dim (\pi) \operatorname{Tr} (\pi (h) u ^ {*} u \pi (h) ^ {- 1}) = \dim (\pi) \operatorname{Tr} (u ^ {*} u) = \| u \| _ {2} ^ {2} \end{array}
\]

for all \((g,h)\in \mathrm{G}\times \mathrm{G}\) and every \(u\in \operatorname {End}(\operatorname {E}_{\pi})\) .

The map \(\Psi\) defined by \(f\mapsto\pi(f)\) is a \((G\times G)\) -morphism from \(\varrho_{G}(\overline{\pi})\) into \(\operatorname{End}(E_{\pi})\) since

\[
\pi (\varrho_ {\mathrm{G}} (g, h) f) = \int_ {\mathrm{G}} f (g ^ {- 1} x h) \pi (x) d \mu (x) = \pi (g) \pi (f) \pi (h ^ {- 1})
\]

for all \((g,h)\in\mathrm{G}\times\mathrm{G}\) and every \(f\in\varrho_{\mathrm{G}}(\overline{\pi})\) . Since \(\varrho_{\mathrm{G}}(\overline{\pi})\) is an irreducible representation (th. 1, \(a\) ) of V, p. 462), there exists \(\lambda\in\mathbf{C}^{*}\) such that the map \(\lambda\Psi\) is zero or isometric (cor. 5, \(a\) ) of V, p. 388). Let \(f=\dim(\pi)\overline{\chi}_{\pi}\in\varrho_{\mathrm{G}}(\overline{\pi})\) . We obtain \(\pi(f)=1_{\mathrm{E}_{\pi}}\) (cor. 2 of V, p. 465), whence \(\|\pi(f)\|_{2}=\dim(\pi)=\|f\|\) (cor. 3 of V, p. 459). Consequently, the map \(\Psi\) is an isometry. Since \(\varrho_{\mathrm{G}}(\overline{\pi})\) and \(\operatorname{End}(\mathrm{E}_{\pi})\) have the same dimension, \(\Psi\) is an isometric isomorphism.

We denote by \(\mathrm{F}^{2}(\widehat{\mathrm{G}})\) the Hilbert sum of Hilbert spaces \(\mathrm{End}(\mathrm{E}_{\pi})\) . The norm of an element \(x \in \mathrm{F}^{2}(\widehat{\mathrm{G}})\) is still denoted \(\|x\|_{2}\) .

In LIE, IX, p. 79, this space is denoted \(\mathrm{L}^{2}(\widehat{\mathrm{G}})\) , a notation we prefer to avoid here so as not to create confusion with the space \(\ell^{2}(\widehat{\mathrm{G}})\) . Let \((u_{\pi})_{\pi\in\widehat{\mathrm{G}}}\) be an element of \(\mathrm{F}^{2}(\widehat{\mathrm{G}})\) . Since

\[
\sum_ {\pi \in \widehat {G}} \| u _ {\pi} \| _ {2} ^ {2} = \sum_ {\pi \in \widehat {G}} \dim (\pi) \operatorname{Tr} (u _ {\pi} ^ {*} u _ {\pi}) <   + \infty
\]

and since \(\|u_{\pi}\|^{2}\leqslant\mathrm{Tr}(u_{\pi}^{*}u_{\pi})\) (cf. EVT, V, p. 52, formula (33)), the norm in \(\mathcal{L}(\mathrm{E}_{\pi})\) of the endomorphism \(u_{\pi}\) tends to 0 at infinity. One can therefore identify \(\mathrm{F}^{2}(\widehat{\mathrm{G}})\) with a subspace of \(\mathrm{F}_{0}(\widehat{\mathrm{G}})\) .

Let \(\pi \in \widehat{\mathbf{G}}\) and \(u \in \mathrm{End}(\mathrm{E}_{\pi})\) . We denote by \(\mathcal{F}_{\pi}(u)\) the function on G defined by \(\mathcal{F}_{\pi}(u)(g) = \langle \pi(g) | u \rangle = \dim(\pi) \mathrm{Tr}(\pi(g)^{*} u)\) for all \(g \in \mathbf{G}\) . It is a continuous function on G. If G is a compact real Lie group, then the function \(\mathcal{F}_{\pi}(u)\) is analytic on G (LIE, III, § 8, no. 1, th. 1).

Since G is compact, one may identify \(L^{2}(G)\) with a subspace of \(L^{1}(G)\) .

THEOREM 2. --- The Fourier cotransformation of G induces, by passing to subspaces, an isometric isomorphism of L²(G) onto F²(Ĝ). Its inverse \(\mathcal{F}_{\mathrm{G}}\) associates to an element \((u_{\pi})_{\pi \in \widehat{\mathrm{G}}}\) of F²(Ĝ) the sum of the series

\[
\sum_ {\pi \in \widehat {\mathrm{G}}} \mathcal {F} _ {\pi} (u _ {\pi}),
\]

which converges in \(\mathrm{L}^2 (\mathrm{G})\)

By the Peter–Weyl theorem (Th. 1 of V, p. 462), the Hilbert space \(L^{2}(G)\) is the Hilbert sum of the spaces \(\varrho_{\mathrm{G}}(\overline{\pi})\) for \(\pi \in \widehat{G}\) . For every \(\pi \in \widehat{G}\) , the linear map \(f \mapsto \pi(f)\) from \(\varrho_{\mathrm{G}}(\overline{\pi})\) into \(\operatorname{End}(\operatorname{E}_{\pi})\) is an isometric isomorphism (Lemma 4). Consequently, the restriction to \(L^{2}(G)\) of the Fourier cotransformation defines an isometric isomorphism of \(L^{2}(G)\) onto \(F^{2}(\widehat{G})\) .

Let \(f \in \mathrm{L}^2(\mathrm{G})\) . Let \(\pi \in \widehat{\mathrm{G}}\) and let \(f_{\overline{\pi}} \in \mathcal{C}(\mathrm{G})\) be the orthogonal projection of \(f\) onto \(\varrho_{\mathrm{G}}(\overline{\pi})\) (th. 1 of V, p. 462). Since this space is the \(\overline{\pi}\)-isotypic component of \(\pmb{\delta}_{\mathrm{G}}\) (cor. 4 of V, p. 463), we have \(f_{\overline{\pi}} = \dim(\pi)\pmb{\delta}_{\mathrm{G}}(\chi_{\pi})(f)\) by cor. 2 of V, p. 465. For every \(x \in \mathrm{G}\) , it follows that

\[
\begin{array}{r l} & {f _ {\overline {{\pi}}} (x) = \dim (\pi) \int_ {\mathrm{G}} \chi_ {\pi} (g) (\pmb {\delta} _ {\mathrm{G}} (g) f) (x) d \mu (g)} \\ & {\qquad = \dim (\pi) \int_ {\mathrm{G}} \chi_ {\pi} (g) f (x g) d \mu (g)} \\ & {\qquad = \dim (\pi) \int_ {\mathrm{G}} \chi_ {\pi} (x ^ {- 1} y) f (y) d \mu (y)} \\ & {\qquad = \dim (\pi) \mathrm{Tr} (\pi (x ^ {- 1}) \pi (f)) = \langle \pi (x) | \pi (f) \rangle ,} \end{array}
\]

that is \(f_{\overline{\pi}} = \mathcal{F}_{\pi}(\pi(f))\) .

Since \(f\) is the sum in \(\mathrm{L}^2 (\mathrm{G})\) of the family \((f_{\overline{\pi}})\) , by the Peter–Weyl theorem, we obtain

\[
f = \sum_ {\pi \in \widehat {\mathrm{G}}} \mathcal {F} _ {\pi} (\pi (f))
\]

where the series converges in \(L^{2}(G)\) . This proves the theorem.

DEFINITION 2. --- The isometric isomorphism \(\mathcal{F}_{\mathrm{G}}\) of \(\mathrm{F}^2 (\widehat{\mathrm{G}})\) onto \(\mathrm{L}^2 (\mathrm{G})\) , which is the inverse of the isomorphism induced by the Fourier cotransformation, is called the Fourier transform of G. The image of an element of \(\mathrm{F}^2 (\widehat{\mathrm{G}})\) is called its Fourier transform.

Remarks. --- 1) The Fourier transform of \((u_{\pi})_{\pi \in \widehat{\mathrm{G}}} \in \mathrm{F}^2(\widehat{\mathrm{G}})\) is therefore the class in \(\mathrm{L}^2(\mathrm{G})\) of the series

\[
g \mapsto \sum_ {\pi \in \widehat {\mathrm{G}}} \langle \pi (g) | u _ {\pi} \rangle = \sum_ {\pi \in \widehat {\mathrm{G}}} \dim (\pi) \operatorname{Tr} (u _ {\pi} \circ \pi (g) ^ {- 1}).
\]

\begin{enumerate}
\def\labelenumi{\arabic{enumi})}
\setcounter{enumi}{1}
\tightlist
\item
  Let \(f \in \mathrm{L}^2(\mathrm{G})\) . We then have the Plancherel formula
\end{enumerate}

\[
\| f \| ^ {2} = \| \overline {{\mathcal {F}}} _ {\mathrm{G}} (f) \| ^ {2} = \sum_ {\pi \in \widehat {\mathrm{G}}} \| \pi (f) \| ^ {2}.
\]

Moreover, by Theorem 2, \(f\) is the sum in \(\mathrm{L}^2 (\mathrm{G})\) of the family \((f_{\pi})_{\pi \in \widehat{\mathrm{G}}}\) where

\[
\begin{array}{l} f _ {\pi} (x) = \mathcal {F} _ {\pi} (\pi (f)) (x) \\ \qquad = \langle \pi (x) | \pi (f) \rangle = \dim (\pi) \int_ {\mathrm{G}} f (g) \chi_ {\pi} (x ^ {- 1} g) d \mu (g) \end{array}
\]

for all \(x \in G\) and every \(\pi \in \widehat{G}\) .

Suppose that G is commutative. Since the irreducible representations of G are of dimension 1 (cor. 7 of V, p. 390) and the dual group \(\widehat{\mathbf{G}}\) is discrete (prop. 18 of II, p. 233), the algebras \(\mathrm{F}_b(\widehat{\mathbf{G}})\) and \(\mathrm{F}_0(\widehat{\mathbf{G}})\) are identified, respectively, with the algebra \(\mathcal{C}_b(\widehat{\mathbf{G}})\) of continuous bounded functions on \(\widehat{\mathbf{G}}\) and the algebra \(\mathcal{C}_0(\widehat{\mathbf{G}})\) of continuous functions tending to 0 at infinity on \(\widehat{\mathbf{G}}\) . Since G is compact, the Haar measure on \(\widehat{\mathbf{G}}\) dual to \(\mu\) is the counting measure \(\widehat{\mu}\) (prop. 18 of II, p. 233). The Hilbert sum \(\mathrm{F}^2 (\widehat{\mathrm{G}})\) of the spaces \(\operatorname{End}(\mathrm{E}_{\pi})\) for \(\pi \in \widehat{\mathrm{G}}\) is identified with the Hilbert space \(\mathrm{L}^2 (\widehat{\mathrm{G}},\widehat{\mu})\) .

Let \(\nu \in \mathcal{M}^1 (\mathrm{G})\) . Then, for every unitary character \(\chi \in \widehat{\mathrm{G}}\) , we have

\[
\chi (\nu) = \int_ {\mathrm{G}} \chi \nu
\]

which proves that the Fourier cotransformation defined above coincides with the Fourier cotransformation of G defined in II, p. 206.

Every \(f \in \mathcal{L}^2(\mathrm{G})\) is integrable and the formula \(\| \overline{\mathcal{F}}_{\mathrm{G}}(f)\| _2 = \| f\|\) is the Plancherel formula (II, p. 215, theorem 1).

PROPOSITION 10. --- The Fourier cotransformation is an injective morphism of involutive Banach algebras from \(\mathcal{M}^{1}(\mathrm{G})\) into \(\mathrm{F}_{b}(\widehat{\mathrm{G}})\) which sends \(\mathrm{L}^{1}(\mathrm{G})\) into \(\mathrm{F}_{0}(\widehat{\mathrm{G}})\) .

Since G is compact, the space \(\mathcal{C}(\mathrm{G})\) is contained and dense in \(\mathcal{L}^{1}(\mathrm{G})\) and \(\mathcal{L}^{2}(\mathrm{G})\) .

Let \(\nu \in \mathcal{M}^1 (\mathrm{G})\) such that \(\overline{\mathcal{F}}_{\mathrm{G}}(\nu) = 0\) . For every function \(f\) continuous on G, we then have \(\overline{\mathcal{F}}_{\mathrm{G}}(\nu *f) = 0\) . Now \(\nu *f\) belongs to \(\mathcal{C}(\mathrm{G})\) (INT, VIII, p. 152, § 4, n°2, prop. 3). Since, by theorem 2, the Fourier cotransformation is injective on \(\mathrm{L}^2 (\mathrm{G})\) , and a fortiori, on the space \(\mathcal{C}(\mathrm{G})\) , we have \(\nu *f = 0\) for \(f\in \mathcal{C}(\mathrm{G})\) . In particular, it follows that

\[
\int_ {\mathrm{G}} f (x) d \nu (x) = (\nu * \check {f}) (e) = 0
\]

for every function \(f \in \mathcal{C}(\mathrm{G})\) (INT, VIII, loc. cit.), therefore the measure \(\nu\) is zero.

The image of \(\mathcal{C}(\mathrm{G})\) by the Fourier cotransformation is contained in \(\mathrm{F}^2 (\widehat{\mathrm{G}})\) , and a fortiori in \(\mathrm{F}_0(\widehat{\mathrm{G}})\) . Since \(\mathrm{F}_0(\widehat{\mathrm{G}})\) is closed in \(\mathrm{F}_b(\widehat{\mathrm{G}})\) and since \(\mathcal{C}(\mathrm{G})\) is dense in \(\mathrm{L}^1 (\mathrm{G})\) , the image of \(\mathrm{L}^1 (\mathrm{G})\) by the Fourier cotransformation is also contained in \(\mathrm{F}_0(\widehat{\mathrm{G}})\) .

PROPOSITION 11. --- a) Let \(u = (u_{\pi})_{\pi \in \widehat{\mathbf{G}}}\) be an element of \(\mathrm{F}(\widehat{\mathbf{G}})\) . If the family \((\mathcal{F}_{\pi}(u_{\pi}))_{\pi \in \widehat{\mathbf{G}}}\) is uniformly summable in \(\mathcal{C}(\mathbf{G})\) , then its sum \(f\) is a continuous function on \(\mathbf{G}\) whose Fourier cotransform is \(u\) ;

\begin{enumerate}
\def\labelenumi{\alph{enumi})}
\setcounter{enumi}{1}
\tightlist
\item
  Let \(f \in \mathrm{L}^1(\mathrm{G})\) . If the family \((\mathcal{F}_\pi(\pi(f)))_{\pi \in \widehat{\mathrm{G}}}\) is uniformly summable in \(\mathcal{C}(\mathrm{G})\) , then its sum is a continuous function on \(\mathrm{G}\) whose class in \(\mathrm{L}^1(\mathrm{G})\) is equal to \(f\) .
\end{enumerate}

Let us prove assertion \(a\) ). Since G is compact, the sum \(f\) of the family \((\mathcal{F}_{\pi}(u_{\pi}))\) belongs to \(\mathrm{L}^2 (\mathrm{G})\) and the series

\[
\sum_ {\pi \in \widehat {\mathrm{G}}} \mathcal {F} _ {\pi} (u _ {\pi})
\]

converges in \(L^{2}(G)\) to f (INT, IV, p. 127, § 3, n° 3, prop. 4). We therefore have \(\mathcal{F}_{\mathrm{G}}((u_{\pi})_{\pi}) = f\) in \(L^{2}(G)\) , whence \((u_{\pi})_{\pi \in \widehat{\mathrm{G}}} = \overline{\mathcal{F}}_{\mathrm{G}}(f)\) (theorem 2).

Let us prove assertion b). We may apply assertion a) to the family \((\pi(f))_{\pi \in \widehat{\mathrm{G}}}\) . The sum g of the family \((\mathcal{F}_{\pi}(\pi(f)))_{\pi \in \widehat{\mathrm{G}}}\) is a continuous function, and hence belongs to \(L^{2}(G)\) , such that \(\overline{\mathcal{F}}_{\mathrm{G}}(g) = (\pi(f))_{\pi \in \widehat{\mathrm{G}}}\) . Since \(\overline{\mathcal{F}}_{\mathrm{G}}\) is injective on \(L^{2}(G)\) , we have f = g in \(L^{2}(G)\) , hence in \(L^{1}(G)\) .

The algebra \(\mathcal{C}(\widehat{\mathrm{G}})\) of complex-valued functions on \(\widehat{\mathrm{G}}\) is identified with the center of the algebra \(\mathrm{F}(\widehat{\mathrm{G}})\) by the map that associates to \(f\colon \widehat{\mathrm{G}}\to \mathbf{C}\) the central element \((f(\pi)1_{\mathrm{E}_{\pi}})_{\pi \in \widehat{\mathrm{G}}}\) of \(\mathrm{F}(\widehat{\mathrm{G}})\) . Denote by \(\mathcal{M}^1 (\mathrm{G}^\sharp)\) the space of central bounded measures on G. It is a closed subspace of the Banach algebra \(\mathcal{M}^1 (\mathrm{G})\) . For every measure \(\nu \in \mathcal{M}^1 (\mathrm{G}^\sharp)\) and every representation \(\pi \in \widehat{\mathrm{G}}\) , we have \(\pi (\nu)\in \mathrm{End}_{\mathrm{G}}(\mathrm{E}_{\pi})\) hence \(\pi (\nu)\) is a scalar multiple of the identity endomorphism of \(\mathrm{E}_{\pi}\) by Schur's lemma (prop. 6 of V, p. 386). The restriction of the Fourier cotransformation to \(\mathcal{M}^1 (\mathrm{G}^\sharp)\) is therefore identified with a unital morphism of involutive Banach algebras from \(\mathcal{M}^1 (\mathrm{G}^\sharp)\) into \(\mathcal{C}(\widehat{\mathrm{G}})\) . This morphism is injective; the image of \(\mathrm{L}^1 (\mathrm{G}^\sharp)\) is contained in \(\mathcal{C}_b(\widehat{\mathrm{G}})\) and its restriction to \(\mathrm{L}^2 (\mathrm{G}^\sharp)\) is an isometric isomorphism onto \(\mathrm{L}^2 (\widehat{\mathrm{G}},\widehat{\mu})\) .

